\chapter{Infinite Series}

A team trains a classifier on one million examples. The objective function is the sum of
the per-example losses, and training stops automatically when the increase in the total
falls below $10^{-7}$ per epoch. At the fortieth epoch the process stops by itself. The
model remains bad: each new example contributes only about $10^{-6}$ to the total, so the
total hardly moves, while that total itself is not necessarily heading toward a finite
value.

The suspicion can be tested with the simplest series. The sum
$H_n = 1 + \tfrac{1}{2} + \cdots + \tfrac{1}{n}$ at $n=10^6$ equals
$14.3927267229$. Adding another million terms raises it only to
$15.0858736534$; the addition is $0.6931469306$, almost exactly $\ln 2 = 0.6931471806$.
Each time the number of terms doubles, the harmonic series rises by about $\ln 2$, without
bound: the estimate $\ln n + \gamma$ gives $H_{10^9} \approx 21.3004815$. A stopping rule
based on ``the total has stopped moving'' therefore cannot tell convergence from
divergence; what tells them apart is the rate at which the terms shrink.

Chapter 4 already supplied limits, the Cauchy criterion, monotone convergence, and the
limit superior. Series add no new idea, because a partial sum is merely a sequence. But
partial sums rarely have a closed formula, and computing $s_{10^6}$ does not tell us
whether $s_{10^7}$ is still moving. This chapter turns that question into one that can be
answered by inspecting $a_n$ alone.

\begin{untukanak}
Imagine two jars of coins. In the first jar you save under a strange rule:
one coin on the first day, half a coin on the second, a quarter on the third, and so on,
always half of the day before. You save every day, never taking a day off, for as long as
you like. How full is that jar after a hundred years? The answer is surprising: never more
than two coins. Every addition does shrink, and the sum stops at a number you can name
exactly.

In the second jar the rule is different. The coin on the first day is large, then smaller,
and smaller still. After a month the jar holds only a small heap. After a year it is a
little thicker. You might conclude that this jar too will stop at some number, because the
additions grow thinner. It does not. If you save long enough, the contents will pass any
number you name, even though climbing by a single unit takes years. Two jars, two rules
that look alike, two opposite fates.

The question: without waiting for years, how do we know which jar stops and which does
not? Is there a sign in the coins themselves, rather than in the final result? This
chapter answers it, and the answer is used every day by people who write programs.
\end{untukanak}

\section{Series, partial sums, and series we can compute directly}
\tingkat{basic}

A series\index{series} is a sum that never finishes. Its definition therefore does not use
the word ``sum'' directly, but rather the sequence of running totals\index{partial sums}.
We read the result after some number of terms, then ask where that reading is heading.

\begin{definisi}
Given a sequence $(a_n)$ in $\mathbb{R}$. The $n$-th partial sum is
\[
  s_n = a_1 + a_2 + \cdots + a_n = \sum_{k=1}^{n} a_k .
\]
The symbol $\sum_{n=1}^{\infty} a_n$ is called a series, and that series \emph{converges}
to $s$ when the sequence $(s_n)$ converges to $s$; in that case we write
$\sum_{n=1}^{\infty} a_n = s$ and call $s$ the sum of the series. If $(s_n)$ diverges, the
series is called \emph{divergent}.
\end{definisi}

\kamus{series}{deret}{a sum that never finishes, whose value is read from the running total after some number of terms.}

Note that $s_n - s_{n-1} = a_n$. So a series and a sequence are the same object; only the
question changes: for a sequence the formula for $s_n$ is known, for a series we hold only
$a_n$.

The Cauchy criterion from Chapter 4 can be translated directly, and its translation is the
most basic tool in this chapter.

\begin{teorema}[Cauchy criterion for series]
$\sum a_n$ converges if and only if for every $\varepsilon > 0$ there is an integer $N$
such that
\[
  m \ge n \ge N \quad \Longrightarrow \quad
  \left| \sum_{k=n}^{m} a_k \right| < \varepsilon .
\]
\end{teorema}

\noindent\textbf{Proof.} Since $\sum_{k=n}^{m} a_k = s_m - s_{n-1}$, the condition above is
nothing but the Cauchy condition for the sequence $(s_n)$. By Chapter 4, a real sequence
converges if and only if it is a Cauchy sequence.

This criterion speaks about a cut-off tail of the series, not about a single term. Taking
$m = n$ gives a necessary condition that is very cheap to check.

\begin{akibat}[The $n$-th term test]
If $\sum a_n$ converges, then $a_n \to 0$.
\end{akibat}

\noindent\textbf{Proof.} Take $m=n$ in the Cauchy criterion: for every $\varepsilon>0$
there is $N$ such that $n \ge N$ implies $|a_n| < \varepsilon$. That is exactly the
definition of $a_n \to 0$.

This test only filters. It says that a term that does not tend to zero certainly wrecks
convergence, and it says nothing about terms that do tend to zero. Two quick examples for
the easy side: $\sum \frac{n}{n+1}$ diverges because $\frac{n}{n+1} \to 1$, and
$\sum (-1)^n$ diverges because its terms do not tend to zero. The hard side is the harmonic
series: $\frac{1}{n} \to 0$ but $\sum \frac{1}{n}$ diverges. We prove that in the second
section.

\begin{bacakeras}
Two voices. Read them in turn, then swap roles.

\textbf{Voice A:} I shrink. Every day I am smaller than yesterday.

\textbf{Voice B:} Good. But shrinking does not yet mean stopping.

\textbf{A:} I shrink by half. Half of half. Half again.

\textbf{B:} Then you have a fence. Your sum will stop.

\textbf{A:} I shrink, but only a little. One third. One quarter. One fifth.

\textbf{B:} Careful. Thin does not mean used up.

\textbf{A:} However small my additions, I keep adding.

\textbf{B:} And however long it takes, your sum keeps rising.

Repeat three times, faster each time: shrinking is the business of terms, stopping is the
business of sums. Memorize: what decides is not the term that shrinks, but the remainder
that never runs out.
\end{bacakeras}

The first series that can be computed directly is the geometric series\index{geometric
series}. Each of its terms is obtained by multiplying the previous term by a fixed number.

\begin{teorema}[Geometric series]
For $|r| < 1$,
\[
  \sum_{k=0}^{\infty} r^k = \frac{1}{1-r},
\]
and for $|r| \ge 1$ the series diverges.
\end{teorema}

\noindent\textbf{Proof.} For $r \neq 1$ the identity
\[
  (1-r)\left(1 + r + \cdots + r^n\right) = 1 - r^{n+1},
\]
holds, so $s_n = \frac{1-r^{n+1}}{1-r}$. If $|r|<1$ then $r^{n+1} \to 0$ (Chapter 4,
since $|r|^{n} \to 0$ for $0 \le |r| < 1$), hence $s_n \to \frac{1}{1-r}$. If
$|r| \ge 1$ then $|r|^k \ge 1$ for all $k$, so $s_n$ cannot converge: for $r=1$ clearly
$s_n = n+1$, and for $|r|\ge 1$ with $r\neq 1$ the term $r^k$ does not tend to zero.

With $r = \tfrac{1}{2}$ we get $1 + \tfrac{1}{2} + \tfrac{1}{4} + \cdots = 2$, and with
$r = -\tfrac{1}{2}$ we get $1 - \tfrac{1}{2} + \tfrac{1}{4} - \cdots = \tfrac{2}{3}$. The
second is interesting: its terms alternate in sign and the sum is smaller than twice the
first term, because opposite signs wipe each other out.

The case $r=\tfrac{1}{2}$ has a picture that explains itself: cut the unit square in half,
cut the remainder again, and keep going. Figure~\ref{fig:potong-geometri} shows the pieces
$\tfrac12$, $\tfrac14$, $\tfrac18$ and the remainder $2^{-5} = 0.03125$ after five cuts,
with used area $1 - 2^{-5} = 0.96875$. That remainder shrinks without vanishing; that is
the geometric face of the sentence ``the sum stops''.

\begin{figure}[htbp]
\centering
\begin{tikzpicture}[scale=6.0]
  \draw[fill=layerteal!45,draw=ink] (0.5,0) rectangle (1,1);
  \draw[fill=layerindigo!40,draw=ink] (0,0.5) rectangle (0.5,1);
  \draw[fill=layerorange!50,draw=ink] (0.25,0) rectangle (0.5,0.5);
  \draw[fill=layerjade!45,draw=ink] (0,0.25) rectangle (0.25,0.5);
  \draw[fill=layerplum!40,draw=ink] (0.125,0) rectangle (0.25,0.25);
  \draw[fill=white,draw=ink,densely dashed] (0,0) rectangle (0.125,0.25);
  \draw[ink,thick] (0,0) rectangle (1,1);
  \node[font=\small] at (0.75,0.5) {$\tfrac{1}{2}$};
  \node[font=\small] at (0.25,0.75) {$\tfrac{1}{4}$};
  \node[font=\small] at (0.375,0.25) {$\tfrac{1}{8}$};
  \node[font=\small] at (0.125,0.375) {$\tfrac{1}{16}$};
  \draw[ink,thin,-{Stealth}] (0.06,0.10) -- (-0.30,-0.05);
  \node[font=\scriptsize,anchor=east] at (-0.30,-0.05) {remainder $\tfrac{1}{32}$};
\end{tikzpicture}
\caption{The geometric series of one half as pieces of the unit square. Pieces of size
$\tfrac12$, $\tfrac14$, $\tfrac18$, $\tfrac1{16}$ are laid down in order, and the
remainder still uncovered after four pieces is $\tfrac1{16}$; after five pieces the
remainder is $2^{-5} = 0.03125$ and the used area $1 - 2^{-5} = 0.96875$. The remainder
always shrinks, it never vanishes.}
\label{fig:potong-geometri}
\end{figure}

\begin{tangan}
Materials: one sheet of square paper, a pencil, a ruler.

\begin{enumerate}
\item Fold the paper in half, open it, shade one half. That is the first piece.
\item Fold the unshaded part in half again, open it, shade half of it.
\item Repeat three more times, always only on the part not yet shaded.
\item Check the white part that remains: after five folds only one thirty-second of the
      original paper is left.
\end{enumerate}

What you should see: every new shading is always half of the previous one, and the white
part never vanishes, it only shrinks toward one corner. If the paper is thin enough to fold
ten times, the white part is only $\tfrac{1}{1024} = 0.0009765625$; by the computation it
tends to zero, so the shading finally covers the whole sheet. Now change the rule: shade
$\tfrac12$ of the paper, then $\tfrac13$, then $\tfrac14$. The folds can no longer follow,
and the paper runs out before the series stops.
\end{tangan}

The second family that can be computed directly is the telescoping series\index{telescoping
series}, that is, a series whose terms can be written as the difference of two consecutive
terms of another sequence, so that almost all subtractions cancel each other.

\begin{contoh}
For $n \ge 1$,
\[
  \frac{1}{n(n+1)} = \frac{1}{n} - \frac{1}{n+1}.
\]
Hence
\[
  s_N = \sum_{n=1}^{N} \left(\frac{1}{n} - \frac{1}{n+1}\right)
      = 1 - \frac{1}{N+1},
\]
because all the intermediate terms cancel and only the first and last terms remain. So
$\sum_{n=1}^{\infty} \frac{1}{n(n+1)} = 1$. At $N=100$ that formula gives
$s_{100} = \tfrac{100}{101} = 0.9900990099$, and the remaining tail is
$\tfrac{1}{101} = 0.0099009901$; the two add up to exactly $1$.
\end{contoh}

Notice that in the telescoping example we need no test at all. A closed formula for $s_N$
gives the limit directly. This is what rarely happens. For almost every series, including
the harmonic series and $\sum 1/n^2$, no such formula exists, and we must decide
convergence without ever knowing the sum exactly.

Before turning to the tests, we record that series behave linearly, just as limits of
sequences do.

\begin{proposisi}
If $\sum a_n = A$ and $\sum b_n = B$ converge, then for every $c, d \in \mathbb{R}$ the
series $\sum (c a_n + d b_n)$ converges with sum $cA + dB$. Adding or discarding finitely
many terms does not change convergence, it only changes the sum.
\end{proposisi}

\noindent\textbf{Proof.} The partial sums of the combined series are
$c s_n^{(a)} + d s_n^{(b)}$, and by the algebraic properties of limits in Chapter 4 that
sequence converges to $cA + dB$. For the second statement, suppose terms are discarded
starting at index $M$. The partial sums of the new series are $s_{N} - s_{M-1}$ for some
$N$, that is, the difference of a convergent sequence and a constant, so the convergence is
the same.

As an application, $\sum_{n=1}^{\infty} \left(\frac{1}{2^n} + \frac{1}{n(n+1)}\right) = 1
+ 1 = 2$, because the first part is a geometric series with $r=\tfrac{1}{2}$ starting at
$n=1$.

\section{Series of nonnegative terms: comparison and the $p$-series}
\tingkat{intermediate}

When all terms are nonnegative, the partial sums can only rise or stay put. The question
of convergence then becomes a far simpler one: do the partial sums stop below some fence?

\begin{teorema}
If $a_n \ge 0$ for all $n$, then $\sum a_n$ converges if and only if its sequence of
partial sums is bounded.
\end{teorema}

\noindent\textbf{Proof.} With $a_n \ge 0$ we have $s_{n+1} = s_n + a_{n+1} \ge s_n$, so
$(s_n)$ is nondecreasing. By the monotone convergence theorem in Chapter 4, a monotone
sequence bounded above converges, and a monotone sequence not bounded above diverges to
$+\infty$.

This theorem looks simple, but this whole section is only a variation on it. All we need is
a way to compare an unknown series with a known one, and the stock of known series is still
thin: the geometric series.

\begin{teorema}[Comparison test]
Suppose $0 \le a_n \le c_n$ for $n \ge N_0$.
\begin{enumerate}
\item If $\sum c_n$ converges, then $\sum a_n$ converges.
\item If $a_n \ge d_n \ge 0$ for $n \ge N_0$ and $\sum d_n$ diverges, then $\sum a_n$
      diverges.
\end{enumerate}
\end{teorema}

\noindent\textbf{Proof.} For (1), take $\varepsilon>0$. By the Cauchy criterion there is
$N \ge N_0$ such that $m \ge n \ge N$ implies $\sum_{k=n}^{m} c_k < \varepsilon$. Since
every term is nonnegative and $a_k \le c_k$,
\[
  \left| \sum_{k=n}^{m} a_k \right| = \sum_{k=n}^{m} a_k \le \sum_{k=n}^{m} c_k
  < \varepsilon ,
\]
so $\sum a_n$ satisfies the Cauchy criterion. For (2), if $\sum a_n$ converged then by (1)
and $d_n \le a_n$ the series $\sum d_n$ would converge as well, contradicting the
hypothesis.

The comparison test\index{comparison test} demands that we pick the right comparison. For
terms in the form of rational fractions, the natural comparison is the largest term in the
numerator divided by the largest term in the denominator. For instance
$\frac{1}{n^2 + 3n} \le \frac{1}{n^2}$, so $\sum \frac{1}{n^2+3n}$ converges provided
$\sum \frac{1}{n^2}$ converges. Hence the whole question moves to a single family of
series.

That family is the $p$-series\index{p-series}, and the tool for handling it is grouping the
terms into blocks whose size is a power of two. This grouping idea is used repeatedly,
including in the Cauchy condensation test\index{Cauchy condensation} below.

\begin{lemma}[Cauchy condensation]
Suppose $a_1 \ge a_2 \ge a_3 \ge \cdots \ge 0$. Then $\sum_{n=1}^{\infty} a_n$ converges
if and only if $\sum_{k=0}^{\infty} 2^k a_{2^k}$ converges.
\end{lemma}

\noindent\textbf{Proof.} Since the terms are nonnegative, both series converge exactly when
their partial sums are bounded (the first theorem of this section). We compare blocks. For
$k \ge 0$, the block $a_{2^k+1} + \cdots + a_{2^{k+1}}$ contains $2^k$ terms, each no
larger than $a_{2^k}$ and no smaller than $a_{2^{k+1}}$. Thus
\[
  \sum_{n=1}^{2^m} a_n \;\le\; a_1 + \sum_{k=0}^{m-1} 2^k a_{2^k},
  \qquad
  \sum_{n=1}^{2^m} a_n \;\ge\; \frac{1}{2}\sum_{k=1}^{m} 2^k a_{2^k}.
\]
The left inequality shows that the partial sums of $\sum a_n$ are bounded when
$\sum 2^k a_{2^k}$ is bounded. The right inequality shows the converse, because the partial
sums of $\sum 2^k a_{2^k}$ are bounded when $\sum a_n$ is bounded. Both directions are
done.

\begin{teorema}[The $p$-series]
$\sum_{n=1}^{\infty} \frac{1}{n^p}$ converges if $p > 1$ and diverges if $p \le 1$.
\end{teorema}

\noindent\textbf{Proof.} For $p \le 0$ the term $\frac{1}{n^p} = n^{-p}$ does not tend to
zero, so the series diverges by the $n$-th term test. For $p>0$ the sequence
$a_n = n^{-p}$ decreases to zero, so the condensation lemma applies and
\[
  \sum_{k=0}^{\infty} 2^k a_{2^k}
  = \sum_{k=0}^{\infty} 2^k \cdot 2^{-kp}
  = \sum_{k=0}^{\infty} \left(2^{1-p}\right)^{k}.
\]
That is a geometric series with ratio $r = 2^{1-p}$. Since $p>1$ exactly when $1-p<0$,
that is exactly when $r<1$, the geometric series theorem gives: convergence for $p>1$,
divergence for $p \le 1$.

The case $p=1$ is the harmonic series\index{harmonic series}, and its proof is now a single
line: the geometric ratio is $2^{1-1} = 1$, so the condensed series $\sum 1$ diverges. We
can read its fragility straight off the blocks: every block of size $2^k$ contributes at
least $2^k \cdot \frac{1}{2^{k+1}} = \tfrac{1}{2}$. Hence
\[
  H_{2^m} \;\ge\; 1 + \frac{m}{2},
\]
so $H_n \to \infty$, but raising $H$ by one requires doubling $n$ many times over: the
growth rate is logarithmic. This is why the number $14.39$ at the opening of this chapter
looks like a bound although it is not.

Figure~\ref{fig:batang-suku} puts the two kinds of terms side by side. The harmonic bars
stay plump: the tenth term is still $0.1$, whereas the tenth geometric term is
$2^{-10} = 0.0009765625$. Sum of the first ten terms: $H_{10} = 2.928968254$ against
$1 - 2^{-10} = 0.9990234375$. The harmonic term is larger at every index, so the size of
the term misleads; what decides is whether the shrinking leaves a finite tail.

\begin{figure}[htbp]
\centering
\begin{tikzpicture}
\begin{axis}[
  width=0.93\linewidth, height=6.0cm,
  ybar, bar width=3.4pt,
  xlabel={index $n$}, ylabel={size of term $a_n$},
  xmin=0.3, xmax=10.7, ymin=0, ymax=1.12,
  xtick={1,2,3,4,5,6,7,8,9,10},
  scaled y ticks=false,
  tick label style={font=\scriptsize},
  label style={font=\small},
  ymajorgrids, grid style={layergrey!20},
  legend style={at={(0.98,0.96)},anchor=north east,draw=none,fill=none,
                font=\scriptsize},
]
\addplot[fill=layerteal!85,draw=layerteal] coordinates {
  (1,1) (2,0.5) (3,0.3333333) (4,0.25) (5,0.2)
  (6,0.1666667) (7,0.1428571) (8,0.125) (9,0.1111111) (10,0.1)};
\addlegendentry{harmonic $1/n$}
\addplot[fill=layerorange!85,draw=layerorange] coordinates {
  (1,0.5) (2,0.25) (3,0.125) (4,0.0625) (5,0.03125)
  (6,0.015625) (7,0.0078125) (8,0.0039063) (9,0.0019531) (10,0.0009766)};
\addlegendentry{geometric $2^{-n}$}
\end{axis}
\end{tikzpicture}
\caption{Bars of the harmonic terms $1/n$ and the geometric terms $2^{-n}$ for
$n=1,\ldots,10$. The harmonic bars are taller at every index and the tenth term is still
$0.1$, whereas the tenth geometric term is $2^{-10} = 0.0009765625$. Sum of the first ten
terms: $H_{10} = 2.928968254$ against $1 - 2^{-10} = 0.9990234375$.}
\label{fig:batang-suku}
\end{figure}

\kamus{harmonic series}{deret harmonik}{a pile that keeps growing although each addition is thinner than the last, so the total height climbs very slowly without bound.}

The $p$-series is not the final boundary. With the same condensation, Rudin shows that
$\sum \frac{1}{n(\log n)^p}$ converges for $p>1$ and diverges for $p<1$, and that adding
one more logarithm shifts the boundary. Those series differ very little yet lie on opposite
sides: there is no tidy boundary line, only tests that hold on a particular class.

\begin{tantangan}
Take $k(n)$ the integer with $2^{k-1} \le n < 2^k$, then set $a_n = \frac{1}{k(n)}$. So
$a_1 = 1$, $a_2 = a_3 = \tfrac12$, $a_4 = \cdots = a_7 = \tfrac13$, and so on. Show that
$a_n \to 0$ but $\sum a_n$ diverges, then explain why the ratio test gives limit superior
$1$ here and is therefore useless.

\textbf{Hint 1.} Do not add the terms one by one. Add one block, and first count how many
terms that block contains.

\textbf{Hint 2.} The $k$-th block contains $2^{k-1}$ terms, all of them equal to
$\frac{1}{k}$, so its contribution is $\frac{2^{k-1}}{k}$. The first ten blocks already
add up to $118.653968$.

If you are stuck, go back to the section on series of nonnegative terms, where blocks of
power-of-two size are used in the Cauchy condensation lemma.
\end{tantangan}

\begin{untukanak}
Try to imagine stacking coins on a table. The first coin is thick. The second is half as
thick. The third a third, the fourth a quarter, and so on: the millionth coin is very thin,
as thin as one millionth of the first. If you stack a million coins, the height of the
stack is about fourteen times the first coin, not a million times. It looks as though the
stack will stop growing, because the next coin adds almost nothing.

Now the question that trips up many people: how many coins does it take for the stack to
reach twenty times the first coin? The answer is about two hundred and seventy-two million
coins. For twenty-one times, one needs about nine billion coins. Each level demands many
times more coins, but no level is impossible. The pile never stops, it only slows down.

This is why computer programs are easily fooled. If you stop stacking after a million
coins, the stack looks as though it no longer moves, when in fact it is not finished. In
this chapter we will know why slow does not mean stopped, and how to tell the difference
before spending the time to stack.
\end{untukanak}

\section{The ratio test and the root test}
\tingkat{intermediate}

The comparison test demands that we already have a comparison series. For terms containing
factorials or powers of $n$, such a comparison is hard to guess, but the comparison between
consecutive terms is easy to compute. The idea: if every term is at most $\beta$ times the
previous one with $\beta<1$, the series shrinks at least as fast as a geometric series, so
it converges.

\begin{teorema}[Ratio test]
Suppose $a_n \neq 0$ for all sufficiently large $n$.
\begin{enumerate}
\item If $\limsup_{n\to\infty} \left| \frac{a_{n+1}}{a_n} \right| < 1$, then
      $\sum a_n$ converges, and in fact converges absolutely.
\item If there is $n_0$ such that $\left| \frac{a_{n+1}}{a_n} \right| \ge 1$ for all
      $n \ge n_0$, then $\sum a_n$ diverges.
\end{enumerate}
\end{teorema}

\noindent\textbf{Proof.} For (1), write $L = \limsup |a_{n+1}/a_n|$ and choose $\beta$ with
$L < \beta < 1$. By the properties of the limit superior in Chapter 4, only finitely many
$n$ satisfy $|a_{n+1}/a_n| > \beta$; so there is $N$ such that $n \ge N$ implies
$|a_{n+1}| \le \beta |a_n|$. By induction,
\[
  |a_{N+k}| \le \beta^{k} |a_N| \qquad (k \ge 0).
\]
The series $\sum \beta^k |a_N|$ converges because $0<\beta<1$, so by the comparison test
$\sum |a_n|$ converges, and by the absolute convergence theorem in the fifth section of
this chapter, absolute convergence implies convergence.

For (2), the condition $|a_{n+1}| \ge |a_n|$ for $n \ge n_0$ makes $(|a_n|)$ nondecreasing
and nonnegative, so it cannot tend to zero. Hence $a_n$ does not tend to zero, and the
$n$-th term test rules out convergence.

\kamus{ratio test}{uji rasio}{measures whether each term shrinks by at least a fixed fraction of the previous one, so that the series shrinks like a geometric series.}
\index{ratio test}

The bound $1$ in the ratio test is a blind spot. For $1/n$ the ratio is
$\frac{n}{n+1} \to 1$ and the series diverges; for $1/n^2$ the ratio is
$\left(\frac{n}{n+1}\right)^2 \to 1$ and the series converges. So a ratio tending to $1$
gives no information; shrinking rates that differ greatly can produce the same ratio.

The next two series show the ratio test working in opposite directions. For
$a_n = \frac{n!}{10^n}$ we have $\frac{a_{n+1}}{a_n} = \frac{n+1}{10}$, which tends to
$+\infty$; the series diverges. For $a_n = \frac{10^n}{n!}$ we have
$\frac{a_{n+1}}{a_n} = \frac{10}{n+1}$, which tends to $0$; the series converges. The sum
of the second series is exactly $e^{10} = 22026.4657948$, whereas the sum of its first $20$
terms is only $21950.3788494$.

The root test\index{root test} uses a different quantity: the $n$-th root of the $n$-th
term. Intuitively it measures the ``average geometric ratio'' of all the terms, not the
ratio between two consecutive terms.

\begin{teorema}[Root test]
\begin{enumerate}
\item If $\limsup_{n\to\infty} \sqrt[n]{|a_n|} < 1$, then $\sum a_n$ converges, and in
      fact converges absolutely.
\item If $\limsup_{n\to\infty} \sqrt[n]{|a_n|} > 1$, then $\sum a_n$ diverges.
\end{enumerate}
\end{teorema}

\noindent\textbf{Proof.} For (1), choose $\beta$ with
$\limsup \sqrt[n]{|a_n|} < \beta < 1$. Only finitely many $n$ satisfy
$\sqrt[n]{|a_n|} > \beta$, so there is $N$ such that $|a_n| \le \beta^n$ for $n \ge N$. The
comparison test with the geometric series $\sum \beta^n$ closes this case.

For (2), if the limit superior is larger than $1$, then $\sqrt[n]{|a_n|} > 1$, that is
$|a_n| > 1$, for infinitely many $n$. Those terms do not tend to zero, so the series
diverges by the $n$-th term test.

The root test is stronger than the ratio test. Let $a_n = 2^{-n}$ for even $n$ and
$a_n = 3^{-n}$ for odd $n$. Then $\sqrt[n]{a_n}$ equals $\tfrac{1}{2}$ or $\tfrac{1}{3}$,
so its limit superior is $\tfrac{1}{2}$ and the root test declares the series convergent.
The ratio swings: at $n=100$, $\frac{a_{101}}{a_{100}}$ equals about
$3.05 \times 10^{17}$, with limit superior $+\infty$ and limit inferior $0$. The ratio test
concludes nothing exactly when the terms swing on different scales.

The ratio test stays popular because factorials and powers of $n$ cancel in the ratio,
whereas the $n$-th root of a factorial demands Stirling's approximation. Order of use: the
$n$-th term test, comparison with a $p$-series or a geometric series, then the ratio test;
the root test is used when the ratio swings or gives $1$.

Now the tests pile up, and the order of checking becomes a skill of its own.

\begin{resep}
\begin{enumerate}
\item Check $a_n \to 0$. If not, the series diverges; done.
\item Recognize the shape: factorials call for the ratio test, powers of $n$ call for the
      root test, rational fractions call for comparison with a $p$-series.
\item If the test you chose gives the bound $1$, conclude nothing. Move to comparison or
      condensation.
\item If the signs alternate and $|a_n|$ decreases to zero, use Leibniz and report the
      error bound $b_{N+1}$.
\end{enumerate}

Warning on order: do not swap step 1 with step 2. The ratio test on a series whose terms do
not tend to zero can give the bound $1$ and mislead, whereas the $n$-th term test settles
it with no computation.
\end{resep}

\section{Alternating series and the Leibniz test}
\tingkat{intermediate}

Up to here all series have nonnegative terms, and there convergence means the partial sums
stop rising. Once the signs alternate, the partial sums can rise and fall, and cancellation
between signs can produce convergence that the size of the terms does not guarantee.

\begin{teorema}[Leibniz test]
Suppose $(b_n)$ is a nonnegative nonincreasing sequence, that is
$b_1 \ge b_2 \ge \cdots \ge 0$, and $b_n \to 0$. Then
$\sum_{n=1}^{\infty} (-1)^{n+1} b_n$ converges.
\end{teorema}

\noindent\textbf{Proof.} Write $s_N$ for the $N$-th partial sum. Consider the partial sums
with even index and with odd index separately. Since
\[
  s_{2m+2} - s_{2m} = b_{2m+1} - b_{2m+2} \ge 0,
\]
the sequence $(s_{2m})$ is nondecreasing. Since
\[
  s_{2m+1} - s_{2m-1} = -b_{2m} + b_{2m+1} \le 0,
\]
the sequence $(s_{2m-1})$ is nonincreasing. Besides, $s_{2m} \le s_{2m-1}$ and
$s_{2m-1} - s_{2m} = b_{2m} \to 0$. So $(s_{2m})$ is nondecreasing and bounded above by
$s_1 = b_1$, hence converges to some $s$; $(s_{2m-1})$ is nonincreasing and bounded below
by $s_2$, hence converges to some $t$; and $s - t = \lim (s_{2m} - s_{2m-1}) = 0$. Both
subsequences tend to the same value, and since the even and odd subsequences together cover
all indices, the whole sequence $(s_N)$ converges to that $s$.

The Leibniz test\index{Leibniz test} gives more than convergence: an error bound that can
be used without knowing the true sum, and that is what makes alternating series useful for
computing constants.

\begin{teorema}[Leibniz error bound]
Under the hypotheses of the Leibniz test and with $s = \sum (-1)^{n+1} b_n$,
\[
  \left| s - s_N \right| \le b_{N+1}.
\]
Sharper: $s$ lies between $s_N$ and $s_{N+1}$, so the sign of the error $s - s_N$ is the
sign of the first discarded term, namely $(-1)^{N+2}$.
\end{teorema}

\noindent\textbf{Proof.} From the earlier proof, $(s_{2m})$ rises toward $s$ and
$(s_{2m-1})$ falls toward $s$. So for every $N$, the number $s$ lies between $s_N$ and
$s_{N+1}$. Since $|s_{N+1} - s_N| = b_{N+1}$, the distance from $s$ to $s_N$ does not
exceed $b_{N+1}$, and the direction of the shift is set by the sign of
$s_{N+1} - s_N = (-1)^{N+2} b_{N+1}$.

The classic example is the alternating harmonic series,
\[
  1 - \frac{1}{2} + \frac{1}{3} - \frac{1}{4} + \cdots = \ln 2 .
\]
The equality with $\ln 2$ is proved via the Taylor series for $\ln(1+x)$ in Chapter 14;
here we only use its convergence, which the Leibniz test guarantees because
$b_n = \frac{1}{n}$ decreases to zero. The number $\ln 2$ appears again, now as the sum of
a convergent series; the difference $H_{2n} - H_n \to \ln 2$ is another version of the same
computation.

Figure~\ref{fig:tangga-leibniz} draws the partial sums of that series as two staircases:
the odd ones descend from $s_1 = 1$, the even ones ascend from $s_2 = 0.5$, and both
bracket $\ln 2 = 0.6931471806$. At $n = 10$ and $n = 11$ the bracket is
$0.6456349206 < \ln 2 < 0.7365440115$, of width $b_{11} = \tfrac{1}{11} = 0.0909090909$,
exactly as the error bound theorem promises.

\begin{figure}[htbp]
\centering
\begin{tikzpicture}
\begin{axis}[
  width=0.93\linewidth, height=6.2cm,
  xlabel={number of terms $n$}, ylabel={partial sum $s_n$},
  xmin=0.4, xmax=15.0, ymin=0.45, ymax=1.06,
  xtick={1,2,3,4,5,6,7,8,9,10,11,12,13,14},
  tick label style={font=\scriptsize},
  label style={font=\small},
  ymajorgrids, grid style={layergrey!20},
  legend style={at={(0.98,0.04)},anchor=south east,draw=layergrey!40,fill=white,
                fill opacity=0.9,text opacity=1,font=\scriptsize},
]
\addplot[color=layerindigo,mark=*,mark size=1.6pt,thick] coordinates {
  (1,1.0) (3,0.8333333) (5,0.7833333) (7,0.7595238)
  (9,0.7456349) (11,0.7365440) (13,0.7301338)};
\addlegendentry{$n$ odd, descending}
\addplot[color=layerteal,mark=square*,mark size=1.6pt,thick] coordinates {
  (2,0.5) (4,0.5833333) (6,0.6166667) (8,0.6345238)
  (10,0.6456349) (12,0.6532107) (14,0.6587052)};
\addlegendentry{$n$ even, ascending}
\draw[dashed,accent,thick] (axis cs:0.4,0.6931472) -- (axis cs:15.0,0.6931472);
\node[font=\scriptsize,text=accent,fill=white,inner sep=1pt]
  at (axis cs:6.4,0.668) {$\ln 2 = 0.6931471806$};
\end{axis}
\end{tikzpicture}
\caption{Partial sums of the alternating harmonic series. The odd subsequence descends,
the even subsequence ascends, and both bracket the limit line $\ln 2 = 0.6931471806$.
At $n = 10$ and $n = 11$: $s_{10} = 0.6456349206$ and $s_{11} = 0.7365440115$, of width
$b_{11} = \tfrac{1}{11} = 0.0909090909$.}
\label{fig:tangga-leibniz}
\end{figure}

The Leibniz bound is useful and is often estimated too optimistically. With $1000$ terms,
the partial sum is $0.7851481635$ whereas $\pi/4 = 0.7853981634$; the actual error is
$2.5000 \times 10^{-4}$, the promised bound
$b_{1001} = \tfrac{1}{2001} = 4.9975 \times 10^{-4}$. The bound is correct but twice the
real error, because the actual error is about $\tfrac{1}{4N}$ whereas the bound is
$\tfrac{1}{2N}$. For an error below $10^{-6}$ we need $N \ge 500000$, and at that $N$ the
error is $5.0000 \times 10^{-7}$.

\section{Absolute convergence, grouping, and rearrangement}
\tingkat{advanced}

Series that converge because their signs cancel are far more fragile than series that
converge because their terms shrink fast enough. That fragility has a name, and this
chapter owes one explanation of it.

\begin{definisi}
The series $\sum a_n$ \emph{converges absolutely} if $\sum |a_n|$ converges. If
$\sum a_n$ converges but $\sum |a_n|$ diverges, the series is called
\emph{conditionally convergent}\index{conditional convergence}.
\end{definisi}

\kamus{absolute convergence}{konvergensi mutlak}{a series that still converges after all its signs are forced positive, so it does not depend on cancellation.}
\index{absolute convergence}

\begin{teorema}
If $\sum a_n$ converges absolutely, then $\sum a_n$ converges and
$\left| \sum a_n \right| \le \sum |a_n|$.
\end{teorema}

\noindent\textbf{Proof.} For every $n \le m$ we have
$\left| \sum_{k=n}^{m} a_k \right| \le \sum_{k=n}^{m} |a_k|$. Since $\sum |a_n|$
converges, the rightmost expression can be made smaller than $\varepsilon$ by the Cauchy
criterion, so $\sum a_n$ satisfies the Cauchy criterion as well. The absolute value
inequality for the sum is obtained by taking limits in
$\left| s_n \right| \le \sum_{k=1}^{n} |a_k|$.

The converse does not hold: the alternating harmonic series converges conditionally,
because $\sum \frac{1}{n}$ diverges. For such series, the order of summation begins to
matter. Before going there, we first look at an operation that is safe, namely
grouping\index{grouping}.

\begin{proposisi}[Grouping]
Suppose $\sum a_n$ converges to $s$, and $n_1 < n_2 < n_3 < \cdots$ is a sequence of
indices. Then the series whose terms are
\[
  A_k = a_{n_{k-1}+1} + a_{n_{k-1}+2} + \cdots + a_{n_k}
\]
(with $n_0 = 0$) also converges to $s$. The converse fails: a convergent grouping does not
guarantee that the original series converges.
\end{proposisi}

\noindent\textbf{Proof.} The partial sums of the grouped series are
$S_K = A_1 + \cdots + A_K = s_{n_K}$, that is, a subsequence of $(s_n)$. A subsequence of a
convergent sequence tends to the same limit (Chapter 4), so $S_K \to s$.

For the converse, take $1 - 1 + 1 - 1 + \cdots$. This series diverges because its terms do
not tend to zero, but grouping in pairs gives $(1-1) + (1-1) + \cdots = 0 + 0 + \cdots$,
which converges to $0$.

That counterexample matters in practice: parentheses may be inserted into a series already
known to converge, but they may not be used to prove it.

Now we change the order, not merely the grouping: a rearrangement renumbers the same terms
so that each term appears exactly once, but at a different position.

\begin{definisi}
Let $\sigma$ be a bijection from the set of positive integers to itself. The series
$\sum a_{\sigma(n)}$ is called a \emph{rearrangement} of $\sum a_n$.
\end{definisi}

\kamus{rearrangement}{susunan ulang}{swapping the order in which goods arrive without changing the goods, so the running total can differ even though the warehouse contents are the same.}
\index{rearrangement}

For an absolutely convergent series, rearrangement is harmless.

\begin{teorema}
If $\sum a_n$ converges absolutely, then every rearrangement of it converges, and all of
them tend to the same sum.
\end{teorema}

\noindent\textbf{Proof.} Write $s = \sum a_n$ and let $\sum a_{\sigma(n)}$ be a
rearrangement with partial sums $s'_N$. Take $\varepsilon > 0$. By the Cauchy criterion
there is $N_0$ such that $m \ge n \ge N_0$ implies $\sum_{k=n}^{m} |a_k| < \varepsilon$; in
particular $\sum_{k=N_0}^{\infty} |a_k| \le \varepsilon$.

Choose $P$ so that the numbers $1, 2, \ldots, N_0-1$ all appear among
$\sigma(1), \ldots, \sigma(P)$. For $N > P$, the terms $a_1, \ldots, a_{N_0-1}$ cancel each
other in the difference $s_N - s'_N$ (both contain those terms), so all that remains are
terms with index $k \ge N_0$. Hence
\[
  \left| s_N - s'_N \right| \;\le\; \sum_{k \ge N_0} |a_k| \;\le\; \varepsilon
  \qquad (N > P).
\]
Since $s_N \to s$ and $\left| s'_N - s \right| \le \left| s'_N - s_N \right| +
\left| s_N - s \right|$, we get $s'_N \to s$.

For series that converge only conditionally, that conclusion collapses completely. The next
theorem is stated without proof; the proof is long and technical, and can be read in Rudin,
in the chapter on sequences and series of numbers, the section ``Rearrangements''.

\begin{teorema}[Riemann rearrangement theorem, statement]
Suppose $\sum a_n$ converges but does not converge absolutely, and suppose $\alpha \le
\beta$ are two real numbers (possibly $\pm\infty$). Then there is a rearrangement
$\sum a_{\sigma(n)}$ whose partial sums have limit inferior $\alpha$ and limit superior
$\beta$. In particular, for every $c \in \mathbb{R}$ there is a rearrangement that
converges to $c$, and there is a rearrangement that diverges.
\end{teorema}

The idea of the proof is this. If $\sum a_n$ converges conditionally, the series of its
positive terms and the series of absolute values of its negative terms both diverge. We
rearrange by rule: add positive terms until we pass $\beta$, then negative terms until we
drop below $\alpha$, then repeat. Because both series diverge, that step can always be
carried out, and because $a_n \to 0$, the overshoot past the boundary gets smaller. The
result oscillates with limit inferior $\alpha$ and limit superior $\beta$. The complete
construction is in Rudin, Theorem 3.54.

The best-known case is the alternating harmonic series rearranged in the pattern two
positive terms then one negative term,
\[
  1 + \frac{1}{3} - \frac{1}{2} + \frac{1}{5} + \frac{1}{7} - \frac{1}{4} + \cdots
\]
This series converges, but not to $\ln 2$: its sum is $\tfrac{3}{2}\ln 2 = 1.0397207708$,
and the computation is in the next section. Changing the order changes the answer even
though the terms are exactly the same. Addition of real numbers is not commutative without
conditions, and that is a concrete warning for anyone who sums fractions on a computer.

This also explains why the tests in the second and third sections require nonnegative terms.
When the signs alternate, shrinking is not enough: the shrinking must be fast enough to beat
the cancellation.

\begin{untukanak}
Imagine two piles of cards: green marked ``in'', red marked ``out'', each with a number.
You move the cards one at a time, and the balance on the screen changes: green adds, red
subtracts. If you change the order of the moves, does the final balance change?

If the greens and the reds are both large so that they cancel each other, the order
matters. Moving the greens first makes the balance spike before it falls; moving the reds
first makes the balance go negative for a while. The final balance can differ too.

There is one safe situation: if the total of all the numbers, ignoring color, stays bounded,
the order has no effect. That is absolute convergence. In this chapter we will know why, and
why a computer can give different answers for sums that look the same.
\end{untukanak}

\section{Worked examples}
\tingkat{intermediate}

This section works through three computations that form the backbone of the chapter. All the
numbers are checked with Python in the lab section; here the intermediate steps are written
so that they can be repeated with paper and pencil.

\medskip\noindent\textbf{1. The harmonic series at one million terms}

We want $H_{10^6}$ without adding a million terms. The tool is the asymptotic form
\[
  H_n = \ln n + \gamma + \frac{1}{2n} - \frac{1}{12n^2} + O\left(\frac{1}{n^4}\right),
\]
where $\gamma = 0.5772156649$ is the Euler--Mascheroni constant. For $n = 10^6$:
\[
  \ln 10^6 = 13.8155105580, \qquad \ln 10^6 + \gamma = 14.3927262229 .
\]
The first correction term $\frac{1}{2n} = 5 \times 10^{-7}$ gives $14.3927267229$, and the
second correction $-\frac{1}{12n^2} = -8.33 \times 10^{-14}$ is far below the precision we
need. So
\[
  H_{10^6} = 14.3927267228657\ldots
\]
Direct summation with \texttt{math.fsum} gives $14.392726722865724$, in agreement to
$10^{-15}$.

More important is how many terms are needed to reach a given bound. From
$H_n \approx \ln n + \gamma$ we get $n \approx e^{H - \gamma}$. For $H = 20$:
\[
  n \approx e^{20 - 0.5772156649} = e^{19.4227843351} = 2.724006 \times 10^{8}.
\]
A direct search gives $n = 272400600$ as the first index with $H_n > 20$, and $H$ at that
point is $20.0000000016$. For $H = 21$ one needs about $7.4 \times 10^8$ terms; the estimate
at $n = 10^9$ gives $21.3004815$. Each added unit demands multiplying $n$ by
$e \approx 2.718$.

\medskip\noindent\textbf{2. The sum $\sum 1/n^2$ and $\pi^2/6$}

With $N = 10^6$ terms, direct summation gives
\[
  s_{10^6} = 1.6449330668487263, \qquad \frac{\pi^2}{6} = 1.6449340668482264 .
\]
The difference is $9.9999950 \times 10^{-7}$, very close to $\frac{1}{N} = 10^{-6}$; that
makes sense because the tail $\sum_{n>N} n^{-2}$ is about $\frac{1}{N}$. The next order
correction is $-\frac{1}{2N^2} = -5 \times 10^{-13}$, so
\[
  s_{10^6} + \frac{1}{N} - \frac{1}{2N^2} = 1.6449340668482262,
\]
which matches $\pi^2/6$ to $2 \times 10^{-16}$, the precision limit of 64-bit fractions. One
million terms plus one correction already gives machine accuracy.

\medskip\noindent\textbf{3. Rearranging the alternating harmonic series}

Take the pattern two positive terms and one negative: in the $k$-th block we add
$\frac{1}{4k-3} + \frac{1}{4k-1} - \frac{1}{2k}$. After $K$ blocks, all the positive terms
with denominator $1, 3, \ldots, 4K-1$ have been used, and the negative terms with
denominator $2, 4, \ldots, 2K$ have been used as well. So
\[
  S_K = \sum_{j=1}^{2K} \frac{1}{2j-1} - \sum_{j=1}^{K} \frac{1}{2j}
      = \left( H_{4K} - \frac{1}{2}H_{2K} \right) - \frac{1}{2}H_K .
\]
With $H_m \approx \ln m + \gamma$,
\[
  S_K \approx \left(\ln 4K - \tfrac{1}{2}\ln 2K\right) - \tfrac{1}{2}\ln K
  = \ln \frac{4K}{\sqrt{2K}\sqrt{K}} = \ln\left(2\sqrt{2}\right) = \tfrac{3}{2}\ln 2 .
\]
Every term $\gamma$ cancels, and the result does not contain $K$. So
$S_K \to 1.0397207708$, not $\ln 2 = 0.6931471806$. At $K = 10^6$ blocks the numerical
value is $1.0397205208$, an error of $2.5 \times 10^{-7}$ against $\tfrac{3}{2}\ln 2$,
consistent with order $\frac{1}{4K}$ as for the Leibniz series.

\kertas{Numerical summary. $H_{10^6} = 14.3927267229$ and
$\ln 10^6 + \gamma = 14.3927262229$, the difference is $\frac{1}{2n} = 5 \times 10^{-7}$.
$H_n > 20$ is first reached at $n = 272400600$, close to $e^{20-\gamma} = 2.724006 \times
10^8$. The Leibniz series with $1000$ terms gives an error of $2.5000 \times 10^{-4}$
against $\pi/4$, with bound $\frac{1}{2001} = 4.9975 \times 10^{-4}$. The sum
$\sum 1/n^2$ up to $10^6$ terms falls short of $\pi^2/6$ by $9.9999950 \times 10^{-7}$; one
Euler--Maclaurin correction closes it to machine precision. The rearrangement two positive
one negative tends to $\frac{3}{2}\ln 2 = 1.0397207708$, not $\ln 2$.}

\section{Computational lab}
\tingkat{advanced}

The following code runs as is on Python 3.14.4 with NumPy 2.5.2. The first part tests the
accuracy of the harmonic sum and of its asymptotic estimate.

\begin{lab}
\begin{verbatim}
import math
import numpy as np

n = 10**6
a = 1.0 / np.arange(1, n + 1)
vals = a.tolist()
naive = 0.0
for x in vals:
    naive += x

print("np.sum   ", repr(float(np.sum(a))))
print("for-loop ", repr(naive))
print("fsum     ", repr(math.fsum(vals)))
small = a.astype(np.float32)
print("float32  ", repr(float(np.sum(small, dtype=np.float32))))
g = 0.5772156649015328606
print("ln n + g ", repr(math.log(n) + g))
print("1/(2n)   ", 1 / (2 * n))
\end{verbatim}
\end{lab}

The output obtained:
\begin{verbatim}
np.sum    14.392726722865726
for-loop  14.392726722864989
fsum      14.392726722865724
float32   14.392724990844727
ln n + g  14.392726222865807
1/(2n)    5e-07
\end{verbatim}

The first four results should be the same, and three of them differ in the last digit:
rounding error. Adding a million numbers of size $10^{-6}$ into an accumulator of size $14$
throws away low bits. \texttt{float32} loses $1.7 \times 10^{-6}$, larger than the
correction $\frac{1}{12n^2}$ we want to measure; when the error we are after is smaller
than the error of the summation itself, the result means nothing.

The second part checks three numbers from the previous section.

\begin{lab}
\begin{verbatim}
import math
import numpy as np

pi4 = math.pi / 4
s = 0.0
for k in range(1000):
    s += (-1)**k / (2*k + 1)
print("Leibniz 1000 terms", repr(s))
print("error", repr(pi4 - s))
print("bound 1/2001", repr(1/2001))

N = 10**6
m = np.arange(1, N + 1, dtype=np.float64)
S = float(np.sum(1.0 / m**2))
print("sum 1/n^2", repr(S))
print("pi^2/6", repr(math.pi**2/6))
print("corrected", repr(S + 1/N - 1/(2*N**2)))

S = 0.0
for k in range(1, 10**6 + 1):
    S += 1.0/(4*k-3) + 1.0/(4*k-1) - 1.0/(2*k)
print("rearrangement", repr(S))
print("1.5 ln 2", repr(1.5*math.log(2)))
\end{verbatim}
\end{lab}

Its output:
\begin{verbatim}
Leibniz 1000 terms 0.7851481634599485
error 0.00024999993749974525
bound 1/2001 0.0004997501249375312
sum 1/n^2 1.6449330668487263
pi^2/6 1.6449340668482264
corrected 1.6449340668482262
rearrangement 1.0397205208399414
1.5 ln 2 1.0397207708399179
\end{verbatim}

Three things need to be read from that output. First, the Leibniz error
$2.5000 \times 10^{-4}$ is exactly half the bound $4.9975 \times 10^{-4}$; that is no
accident, because the remainder of a monotone alternating series approaches half of the
first discarded term. Second, the Euler--Maclaurin correction closes the gap to
$2 \times 10^{-16}$, one rounding unit at magnitude $1.64$. Third, the rearrangement gives
$1.0397205$, not $0.6931472$, and its error against $\tfrac{3}{2}\ln 2$ is only
$2.5 \times 10^{-7}$.

\begin{untukdosen}
The teaching order I recommend: start from the epoch problem in the opening, then the
definition of partial sums and the two series that can be computed (geometric,
telescoping). The tests come in as a sequence of checks, not a list to memorize: the
$n$-th term, comparison, ratio, root. The $p$-series follows after condensation, which
explains why the harmonic series diverges and why it does so slowly.

Two misconceptions appear most often. First, students conclude convergence from
$a_n \to 0$; the harmonic series is the cure, and it is best shown with a table
$H_{10} = 2.9290$, $H_{100} = 5.1874$, and $H_{1000} = 7.4855$ so that the slowness is
visible. Second, students take ``convergent series'' to mean ``partial sums stop changing'',
when the change per term can be far below machine precision while the series diverges. This
is the right place to mention rounding error.

A class activity that works: split the class into groups, give each group four unnamed
series, and the only task is to decide convergence or divergence together with the test,
in ten minutes. The chosen series must include one trap with ratio $=1$. Discuss groups that
use different tests and are both right. Students are ready to move on when they can explain
why $\sum 1/n^2$ converges while $\sum 1/n$ does not, using only block grouping.
\end{untukdosen}

\section{Summary}
\tingkat{basic}

\begin{ringkasan}
\begin{itemize}
\item A series is a sequence of partial sums; converging means that sequence converges.
\item For nonnegative terms, convergence is equivalent to bounded partial sums. From this
      come the comparison test, condensation, and the classification of $p$-series:
      convergent exactly when $p>1$.
\item The ratio test and the root test measure the speed of shrinking relative to a
      geometric series. Both are blind at the bound $1$, and the root test is stronger than
      the ratio test.
\item The Leibniz test gives convergence for alternating series that decrease to zero,
      along with an error bound $b_{N+1}$ whose sign is known.
\item Absolute convergence is immune to grouping and rearrangement. Conditional
      convergence is not: a Riemann rearrangement can turn the sum into any number, or make
      it diverge.
\item The harmonic series diverges logarithmically: $H_{10^6} \approx 14.39$ and it takes
      $272$ million terms to reach $20$. A stopping rule based on ``the total no longer
      changes'' cannot be trusted.
\end{itemize}
\end{ringkasan}

\section{Exercises}
\tingkat{intermediate}

\begin{latihan}
\begin{enumerate}
\item Compute $s_{20}$ for the geometric series $\sum_{k=0}^{\infty}
      \left(\frac{3}{5}\right)^k$, then compare it with the true sum. What is the error,
      and what error bound can you promise from the formula $\frac{r^{21}}{1-r}$?

\item Prove that $\sum_{n=1}^{\infty} \frac{1}{n(n+2)} = \frac{3}{4}$ by splitting the
      term into partial fractions and showing that the series telescopes.

\item Prove: if $a_n \ge 0$ and $\sum a_n$ converges, then $\sum a_n^2$ converges. Show by
      example that the converse is false.

\item Evaluate this claim: true, false, or not yet enough evidence. ``Because
      $\frac{1}{n} \to 0$ and every term is smaller than the terms of the geometric series
      with ratio $\frac{1}{2}$ for $n$ sufficiently large, the harmonic series converges.''
      State exactly where the reasoning breaks.

\item Decide convergence or divergence, naming the test used:
      (a) $\sum \frac{n^2}{2^n}$, (b) $\sum \frac{2^n}{n^2}$,
      (c) $\sum \frac{n!}{n^n}$, (d) $\sum \frac{1}{n^{1/3}}$.

\item Compute the sum $\sum_{n=1}^{10^6} \frac{1}{n^2}$ with NumPy, then estimate the tail
      with $\frac{1}{N} - \frac{1}{2N^2}$ and compare the result with $\pi^2/6$. How many
      digits after the decimal point are correct before the correction, and after it?

\item The series $\sum \frac{(-1)^{n+1}}{\sqrt{n}}$ converges. Determine the smallest $N$
      such that the Leibniz bound guarantees an error below $10^{-3}$, then compute the
      actual error at that $N$ with Python. Is the bound loose or sharp?

\item Prove that if $\sum a_n$ converges absolutely and $(b_n)$ is bounded, then
      $\sum a_n b_n$ converges absolutely.

\item Rearrange the alternating harmonic series in the pattern three positive terms and one
      negative. Guess the sum with the method of the worked examples section, then check it
      numerically. Is the result still a multiple of $\ln 2$?

\item Prove that if $\sum a_n$ converges conditionally, then the series of its positive
      terms diverges and the series of absolute values of its negative terms diverges as
      well. Explain why this fact is the fuel of the Riemann rearrangement theorem.
\end{enumerate}

\noindent\emph{Hints for the three hardest items.} Item 7: the Leibniz bound uses the
first discarded term, so solve $\frac{1}{\sqrt{N+1}} < 10^{-3}$ and check whether that
inequality is sharp enough. Item 9: after $K$ blocks, the positive terms are used up to
denominator $6K-1$, the negative ones up to $\frac{1}{2K}$; write the partial sums in terms
of $H_{6K}$, $H_{3K}$, and $H_{K}$, then take the limit. Item 10: write
$p_n = \max(a_n, 0)$ and $q_n = \max(-a_n, 0)$; then $a_n = p_n - q_n$ and
$|a_n| = p_n + q_n$, so convergence of one of those series says something about
$\sum |a_n|$.
\end{latihan}

\begin{untukpeneliti}
What is really proved in this chapter are the statements with complete proofs: the Cauchy
criterion for series, the comparison test, Cauchy condensation, the classification of
$p$-series, the ratio test and the root test in limit superior form, the Leibniz test
together with its error bound, and the closure of absolutely convergent series under
rearrangement. I quote the Riemann rearrangement statement without proof. The primary
source is Rudin, \emph{Principles of Mathematical Analysis}, chapter ``Numerical Sequences
and Series'': Theorems 3.22--3.23, 3.25, 3.27--3.28, 3.33--3.34, 3.42--3.43, 3.45--3.46,
and 3.54--3.55. Abbott, \emph{Understanding Analysis}, chapter ``Sequences and Series'',
discusses the same material with emphasis on rearrangement in the opening section of that
chapter.

What is often popularized beyond what is proved is the idea of a threshold between
convergent and divergent series. Rudin shows that no such threshold exists: two series that
differ very little can lie on opposite sides, and the procedure can be continued without
end.

An open question: how reliable is a convergence decision when $a_n$ is available only as a
noisy estimate, such as the sum of losses in machine learning? For such series there is not
yet a tidy theory.

Limits of validity: the whole chapter treats series of real numbers. In other normed spaces
the definition is the same but the tests differ; series of functions require uniform
convergence (Chapter 14). All numbers come from NumPy 2.5.2, 64-bit arithmetic.
\end{untukpeneliti}

\peta{Chapter 5 in Rudin, \emph{Principles of Mathematical Analysis}, chapter ``Numerical
Sequences and Series'': series and partial sums (3.21--3.24), series of nonnegative terms
and comparison (3.25--3.29), the root and ratio tests (3.33--3.37), absolute convergence and
rearrangement (3.45--3.55). In Abbott, \emph{Understanding Analysis}, chapter ``Sequences
and Series'': rearrangement in the opening section of that chapter. This chapter does not
reuse the prose of Chapter 4; from there it refers only to the definition of a limit, the
Cauchy criterion, and monotone convergence.}
